\section{Limitations and intended use}
The platform is a computational assembly with several levels of biological constraint. A single set of LIF parameters, engineering rate conversions, gain corrections, and a simplified plasticity rule cannot reproduce all cellular dynamics. The motor bridge combines different specimens and sexes and transfers rates rather than individual spike timing. No biomechanical body closes the sensorimotor loop. These limits concern the modeled mechanism itself, not merely presentation.

The language interface introduces another source of error. Object descriptions can map to imperfect chemical proxies, a planner can apply the wrong action, and a narrator can state more than the recorded evidence supports. The soft-token reader has a distinct training distribution and an observed failure on held-out taste interaction. Fluent text and animated movement must therefore be interpreted together with explicit inputs and numerical readouts.

The case-study panel characterizes selected mechanisms within this assembly. It does not establish whole-animal validity or biological generalization across individual animals. A second independently modeled fly and a physical social environment are future extensions. The present intended use is reproducible exploration of explicit modeling assumptions and their consequences.
